\documentclass[../../main.tex]{subfiles}
\begin{document}

\begin{center}
  {\fontsize{48pt}{58pt}\selectfont \textbf{1989}\normalsize\textbf{年度}}
\end{center}

\begin{table}[htb]
  \centering
  \caption{出題テーマと難易度}
  \begin{tabular}{|c|p{6.2cm}|c|c|c|}
    \hline
    \textbf{問題番号} & \textbf{テーマ} & \textbf{思考量} & \textbf{計算量} & \textbf{総合} \\
    \hline
    第1問 & 奇数の最大公約数と平方数 & 難 & 普 & 難 \\
    \hline
    第2問 & 放物面と球面で囲まれる体積 & 普 & 普 & 普 \\
    \hline
    第3問 & 二等辺三角形と半平面の共通部分の最大面積 & 普 & 普 & 普 \\
    \hline
    第4問 & 行列による一次変換と三角形の面積 &  &  &  \\
    \hline
    第5問 & カードの得点を比較する確率 & 普 & 難 & 普 \\
    \hline
    第6問 & 分数関数の反復合成と数列の極限 & 普 & 普 & 普 \\
    \hline
  \end{tabular}
\end{table}

\begin{table}[htb]
  \centering
  \caption{解答}
  \renewcommand{\arraystretch}{1.4}
  \begin{tabular}{|c|c|p{9cm}|}
    \hline
    \textbf{問題番号} & \textbf{小問} & \textbf{答え} \\
    \hline
    第1問 & (1) & $2d$，$4d$，$8d$のいずれか（証明問題） \\
    \cline{2-3}
    & (2) & ともに平方数にはならない（証明問題） \\
    \hline
    第2問 & -- & $\pi\left(\dfrac{r^4}{2}-\dfrac{2r^3}{3}+\dfrac{r^2}{4}-\dfrac{1}{96}\right)$ \\
    \hline
    第3問 & -- & 最大値$8\sqrt{2}$ \\
    \hline
    第4問 & -- &  \\
    \hline
    第5問 & -- & A君が勝つ確率の方が大きい \\
    \hline
    第6問 & (1) & $f_n(x)=\dfrac{a^n x}{1+x\sum_{k=0}^{n-1}a^k}$ \\
    \cline{2-3}
    & (2) & $a=1$では$b_n=n$，$0<a<1$では$b_n=a^{1-n}$の一例 \\
    \hline
  \end{tabular}
\end{table}

\paragraph{試験形態}
6問．試験時間は要確認．

\end{document}
